%=============================================================================!
% 1.8  SHORT INTERVALS AND THE TAYLOR SERIES                                  !
%=============================================================================!
\section{Short intervals and the Taylor series}
\label{sec:mo-taylor}

For the ball, the acceleration was constant and the integrals of
\cref{sec:mo-integrating} could be done by hand. For most motions, and for
general systems of interacting atoms, the acceleration changes from
moment to moment and the integrals cannot be done in closed form. A
computer then advances the motion through many short intervals, and on
each interval it needs an approximation to the position a short time
$\tau$ ahead, built from what is known at the start of the interval. The
Taylor series provides that approximation, and tells us how large its
error is.

\subsection*{Matching a polynomial to a function}

Near a time $t$, we look for a polynomial in $\tau$ that matches $x(t +
\tau)$ as closely as possible,
\begin{equation*}
   x(t+\tau) \approx c_0 + c_1\tau + c_2\tau^2 + c_3\tau^3 + \dotsb ,
\end{equation*}
and we choose its coefficients so that the polynomial and the function
have the same value, the same first derivative, the same second
derivative, and so on, at $\tau = 0$. Each condition fixes one
coefficient.
\begin{itemize}
\item Setting $\tau = 0$ leaves only $c_0$, so $c_0 = x(t)$.
\item Differentiating once with respect to $\tau$ gives $c_1 + 2c_2\tau +
   3c_3\tau^2 + \dotsb$, which at $\tau = 0$ is $c_1$. The function's own
   derivative with respect to $\tau$, with $t$ held fixed, is by the chain
   rule with inner function $u = t + \tau$, whose derivative with respect
   to $\tau$ is 1, $\dot{x}(t+\tau)$; repeating this, its $n$th derivative
   is $x^{(n)}(t+\tau)$, the $n$th derivative of $x$ evaluated at
   $t+\tau$. At $\tau = 0$ the first derivative is $\dot{x}(t)$, so matching
   gives $c_1 = \dot{x}(t)$.
\item Differentiating twice gives $2c_2 + 6c_3\tau + \dotsb$, which at
   $\tau = 0$ is $2c_2$, so $c_2 = \ddot{x}(t)/2$.
\item Differentiating three times gives $6c_3 + \dotsb$, so $c_3 =
   \dddot{x}(t)/6$, where $\dddot{x}$ is the third derivative.
\end{itemize}
Each differentiation brings down one more factor from the power of
$\tau$, so the $n$th differentiation of $c_n\tau^n$ produces $n\times(n-1)
\times\dotsb\times 1 = n!$ times $c_n$, a product written $n!$ and read
`$n$ factorial'. After $n$ differentiations the lower powers of $\tau$ have
become zero and the higher ones still contain $\tau$, so at $\tau = 0$ only
$n!\,c_n$ survives, and matching the $n$th derivative $x^{(n)}(t)$ gives
$c_n = x^{(n)}(t)/n!$. Collecting the coefficients gives the Taylor
series,
\begin{equation*}
   x(t+\tau) = \sum_{n=0}^{\infty} \frac{x^{(n)}(t)}{n!}\,\tau^n
   = x(t) + \dot{x}(t)\,\tau + \frac{\ddot{x}(t)}{2}\,\tau^2
   + \frac{\dddot{x}(t)}{6}\,\tau^3 + \dotsb ,
\end{equation*}
with $0! = 1$. For the functions met in this book, polynomials, sines,
cosines and exponentials, the infinite series converges to $x(t+\tau)$:
adding more and more terms brings the sum as close to $x(t+\tau)$ as we
please.
In general, matching derivatives guarantees only that the series stopped
after a finite number of terms has a small error, as described below, and
that is all the book relies on. For a motion, $\dot{x} = v$ and $\ddot{x}
= a$, so
\begin{equation}
   x(t+\tau) = x(t) + v(t)\,\tau + \tfrac12 a(t)\,\tau^2
   + \tfrac16\,\dddot{x}(t)\,\tau^3 + \order{\tau^4} .
   \label{eq:mo-taylor}
\end{equation}

\subsection*{The size of the error}

The symbol $\order{\tau^4}$ stands for everything left out, and it
records how quickly that remainder vanishes as the interval shrinks: a
quantity is $\order{\tau^4}$, read `of order $\tau^4$', if its absolute
value is no larger than some fixed constant times $\lvert\tau\rvert^4$
once $\lvert\tau\rvert$ is small enough. More generally, stopping the
series after the term in $\tau^n$ leaves an error of order $\tau^{n+1}$
for a smooth function, one that can be differentiated as often as we
like. To see why, call the error $D(\tau)$, the function $x(t+\tau)$
minus the polynomial of degree $n$, and take $\tau > 0$. By construction
$D$ and its first $n$ derivatives with respect to $\tau$ vanish at $\tau
= 0$. Since $n + 1$ differentiations remove a polynomial of degree $n$
entirely, the $(n+1)$th derivative of $D$ is $x^{(n+1)}(t+\tau)$. Let $M$
be the greatest value of $\lvert x^{(n+1)}\rvert$ between $t$ and $t +
\tau$. By \cref{eq:mo-integrals} the $n$th derivative of $D$, which
starts at zero, is the integral of the $(n+1)$th from 0 to $\tau$, an
area under a curve lying between $-M$ and $M$, so its size is at most
$M\tau$. Integrating again, the size of the $(n-1)$th derivative is at
most $\int_0^\tau M s\,\dd s = M\tau^2/2$. Each further integration raises
the power by one and divides by the new power, so after $n + 1$
integrations
\begin{equation*}
   \lvert D(\tau)\rvert \le \frac{M\,\lvert\tau\rvert^{n+1}}{(n+1)!} ,
\end{equation*}
and the same argument holds for $\tau < 0$. Halving the interval
therefore divides this bound by $2^{n+1}$, and the error itself by about
that factor.

\begin{toolbox}[Logarithmic axes]
Errors such as these span many powers of ten, so they are plotted on
logarithmic axes, on which the marks $10^{-3}$, $10^{-2}$ and $10^{-1}$
are equally spaced. Equal distances along such an axis stand for equal
ratios, here a factor of ten, rather than equal differences: the
position of a positive number $y$ along the axis is proportional to its
logarithm $\log y$, the power to which 10 must be raised to give $y$, so
that $\log 10^k = k$. Writing $y = 10^a$ and $z = 10^b$, so that $yz =
10^a\times10^b = 10^{a+b}$, gives $\log(yz) = a + b = \log y + \log z$:
the logarithm turns products into sums. Since $(10^a)^p = 10^{pa}$, it
also turns a power into a multiple, $\log(\tau^p) = p\log\tau$. An error
$D$ that behaves as $D = C\tau^p$ therefore obeys
\begin{equation*}
   \log D = \log C + p\log\tau ,
\end{equation*}
which is a straight line when $\log D$ is plotted against $\log\tau$.
Its slope is the power $p$: a line that rises by one power of ten for
every power of ten in $\tau$ has $p = 1$, and one that rises by four has
$p = 4$. The constant $C$ only moves the line up or down. The slope of an
error plot thus gives the power with which the error falls, and two
lines of equal slope differ only by a constant factor.
\end{toolbox}

\Cref{fig:mo-taylor} shows these properties for $x(t) = \sin t$ expanded
about $t_1 = \pi/4$. By \cref{eq:mo-trig-derivs} the derivatives of the
sine cycle through $\sin, \cos, -\sin, -\cos$, so at $\pi/4$ every one of
them is $\pm 1/\sqrt2$, and none is zero, so the leading term
$x^{(n+1)}(t_1)\,\tau^{n+1}/(n+1)!$ of every error is present. The
polynomial of degree 1 is the tangent line,
the polynomial of degree 2 adds the bending of the graph, and each further
degree follows the function over a wider range. Close to $t_1$, the error
of the polynomial of degree $n$ falls as $\tau^{n+1}$, so on the
logarithmic axes of \cref{fig:mo-taylor}(b) it is a straight line of
slope $n + 1$: for the cubic,
halving
$\tau$ from 0.1 to 0.05 divides the error by 16.15, close to $2^4 = 16$.

\begin{figure}[htb]
\centering
\includegraphics{ch01_motion/taylor}
\caption{Taylor polynomials of $x(t) = \sin t$ about $t_1 = \pi/4$. (a)
The polynomials of degrees 1, 2 and 3 (colours as in the key of b) follow
the function ever more closely near $t_1$. (b) The error of the polynomial
of degree $n$ at $t_1 + \tau$, including the constant polynomial of degree
0, with its leading term $\lvert x^{(n+1)}(t_1)\rvert\,\tau^{n+1}/(n+1)!$
(dashed); each error falls as $\tau^{n+1}$.}
\label{fig:mo-taylor}
\end{figure}

\subsection*{Two series used later}

Setting $t = 0$ in the Taylor series and writing $y$ for $\tau$, so that
$x(y) = x(0) + \dot{x}(0)\,y + \ddot{x}(0)\,y^2/2! + \dotsb$, and using
$\sin 0 = 0$ and $\cos 0 = 1$, the
derivatives of $\sin y$ at $y = 0$ are $0, 1, 0, -1, 0, 1, \dots$ and
those of $\cos y$ are $1, 0, -1, 0, 1, \dots$, so
\begin{equation}
   \sin y = y - \frac{y^3}{3!} + \frac{y^5}{5!} - \dotsb, \qquad
   \cos y = 1 - \frac{y^2}{2!} + \frac{y^4}{4!} - \dotsb .
   \label{eq:mo-sin-series}
\end{equation}
For small angles these give the approximations $\sin y \approx y$ and
$\cos y \approx 1 - y^2/2$, which are used throughout the book; the
first agrees with the limit $\sin h/h \to 1$ found in \cref{sec:mo-trig}.

\subsection*{Back to the ball}

For the ball, the Taylor series ends after its third term, because
$\dddot{x} = \dd a/\dd t = 0$ when the acceleration is constant, and
every higher derivative vanishes too. The expansion of
\cref{sec:mo-velocity},
\begin{equation*}
   x(t_1+\tau) = x(t_1) + \bigl(v_0 - g t_1\bigr)\,\tau - \tfrac12 g\,\tau^2 ,
\end{equation*}
was this series in disguise: its coefficients are $v(t_1)$ and
$\tfrac12 a$. For a general motion, the first three terms of
\cref{eq:mo-taylor} amount to \cref{eq:mo-constant} with the acceleration
held at its value at the start of the interval, and their error is of
order $\tau^3$, small when the interval is short. Advancing a motion
through many such intervals is how a computer follows it, and Chapter~9
develops this idea into the methods that molecular dynamics uses.

\notebook{01}{Taylor polynomials of growing order}
